% Generated draft: requires mathematical and source review.
\documentclass[11pt]{article}
\usepackage[margin=25mm]{geometry}
\usepackage{fontspec}
\setmainfont{Latin Modern Roman}
\usepackage{amsmath,amssymb,mathtools,unicode-math}
\setmathfont{Latin Modern Math}
\usepackage{enumitem,needspace,xurl,hyperref}
\setlength{\parindent}{0pt}
\setlength{\parskip}{5pt}
\setlength{\emergencystretch}{3em}
\newcommand{\N}{\mathbb{N}}
\newcommand{\Z}{\mathbb{Z}}
\newcommand{\Q}{\mathbb{Q}}
\newcommand{\R}{\mathbb{R}}
\newcommand{\C}{\mathbb{C}}
\newcommand{\F}{\mathbb{F}}
\newcommand{\A}{\mathbb{A}}
\newcommand{\PP}{\mathbb{P}}
\newcommand{\E}{\mathbb{E}}
\newcommand{\eps}{\varepsilon}
\newcommand{\dd}{\,\mathrm{d}}
\begin{document}
\section*{UnsolvedMath ID: 30005037}
\textbf{Source problem code:} catalogue entry
\subsection{Definitions and mathematical statement}
Let
\[
\Delta=\langle a,b,c,d,e,f,g\mid acb=1,\ add=1,\ afg=1,\ bbf=1,\ cdg=1,\ cef=1,\ eeg=1\rangle.
\]
The supplied record describes \(\Delta\) as torsion-free and as lacking the unique-product property. A group \(G\) is torsion-free if \(g^n=1\), for \(g\in G\) and an integer \(n\geq 1\), implies \(g=1\). It has the unique-product property if, for every pair of finite nonempty subsets \(A,B\subseteq G\), some element of \(AB\) has exactly one representation \(ab\) with \((a,b)\in A\times B\). These descriptions of \(\Delta\) are supplied source assertions, not independently established here.

For a field \(K\), the group algebra \(K[\Delta]\) consists of finite formal sums \(x=\sum_{g\in\Delta}x_g g\), with coefficients \(x_g\in K\). Addition is coefficientwise, and multiplication is the bilinear extension of multiplication in \(\Delta\):
\[
\left(\sum_g x_g g\right)\left(\sum_h y_h h\right)
=\sum_{t\in\Delta}\left(\sum_{gh=t}x_g y_h\right)t.
\]
The question is whether \(\Delta\) satisfies Kaplansky's zero-divisor conjecture over every field; explicitly, whether
\[
\text{for every field }K\text{ and all }x,y\in K[\Delta],
\qquad xy=0\ \Longrightarrow\ x=0\text{ or }y=0.
\]
There is no restriction on the characteristic of \(K\).

\textit{Supplied contextual qualifications.} The original passage says that the unit conjecture remains open in characteristic zero, that Murray subsequently obtained counterexamples using the same group \(P\) in every positive characteristic (its reference [6]), and that the zero-divisor conjecture is known for many groups, including residually torsion-free elementary amenable groups (its reference [5]). It contrasts this with failure of the unit conjecture for the Hantzsche--Wendt group. It also states that failure of the zero-divisor conjecture would yield a counterexample in positive characteristic, and that, for torsion-free \(G\) and a field \(K\) of positive characteristic, the conjecture for \(K[G]\) is equivalent to its unit group having ``no non-trivial torsion.'' These are retained as supplied contextual claims, not independently verified facts. The passage does not define \(P\), identify references [5] and [6] bibliographically, or specify the meaning of ``non-trivial torsion.'' In particular, the latter phrase cannot simply mean absence of all nonidentity finite-order units, since scalar units can have finite order. This unresolved wording is not used to replace the explicit zero-divisor question above.
\subsection{Short English statement}
For the specified seven-generator group \(\Delta\), ask whether the product of two nonzero elements of \(K[\Delta]\) is nonzero for every field \(K\), in every characteristic. The supplied record labels the problem open; this statement-only draft does not independently verify that status or the contextual literature claims.
\subsection{Sources}
\begin{itemize}
\item Local UnsolvedMath registry, supplied complete record with permanent ID \(30005037\), source problem code \texttt{catalogue entry}, in the Oberwolfach Reports Open Problems collection. The record gives proposed year 2022 and names no proposer.
\item Supplied source citation: \emph{Geometric Structures in Group Theory (2023)}, Oberwolfach Reports, DOI 10.4171/owr/2022/11, pp.~517--576. \url{https://doi.org/10.4171/owr/2022/11}. The background attributes the problem to workshop open-problem discussions documented in this report.
\item Supplied literature source: \emph{On zero divisors and units in group rings}. \url{https://arxiv.org/abs/2501.07646}. No authors or further bibliographic details for this item are supplied.
\item Literature assessment dated exactly \texttt{2026-08-22}: the supplied curation labels this question open and says that the general zero-divisor conjecture remains unresolved. This is supplied dated curation, not an independently verified assessment. The background explicitly requests independent verification of time-sensitive status and cited claims.
\item The supplied record labels the statement \texttt{corrected\_verified} and reports that the official Oberwolfach report was checked directly to restore definitions and remove extraction artifacts. This verification claim belongs to the supplied record; no source visit or independent verification is claimed here. The original passage's references [5] and [6] are supplied without their bibliography entries.
\end{itemize}
\end{document}
